% \section{Discussion about network size}\label{sec:appx_lowerbound}
% If we just assume continuity of a permutation equivariant function without any additional smoothness, quantizing the domain to cubes and approximating the function with constants require memorizing $(\text{output dim}) \times (\text{num cubes}) / n!$ real numbers, where $1/n!$ factor is due to permutation equivariance. Thus, one can expect that at least $O(dn(1/\delta)^{dn}/n!)$ parameters are inevitably required. 

% As we state in Lemmas~\ref{lemma:quantize}, \ref{lemma:contextmap}, and \ref{lemma:memorize}, the sum of number of layers required is $O(n(1/\delta)^{dn}/n!)$.
% Note that each transformer block is of constant width, so it requires $O(d)$ parameters.
% This means that the total number of parameters is $O(dn(1/\delta)^{dn}/n!)$, which agrees with our expectation and shows that our construction is optimal in the number of parameters.

% If we compare it with the similar ResNet result, where $X$ is flattened into a vector and we fit the piecewise linear function over the $dn$-dimensional vector, the proof technique in \citep{lin2018resnet} requires $O((1/\delta)^{dn})$ layers, and each layer has $O(dn)$ parameters. So, the total requirement is $O(dn (1/\delta)^{dn})$, which shows that Transformers are better suited for representing permutation equivariant functions.

%%%%%%%%%%%%%%%%%%%%%%%%%%%%%%%%%%%%%%%%%%%%%%%%%%%%%%%


\section{Proof of Claim~\ref{claim:equivariance}}
\label{sec:proof_equivariance}
Suppose $\mX \mP$ was given as input, where $\mP$ is a permutation matrix.
First note that 
\begin{equation*}
(\mW_K^i \mX \mP)^T (\mW_Q^i \mX \mP) = \mP^T (\mW_K^i \mX)^T (\mW_Q^i \mX) \mP    
\end{equation*}
After the softmax operation, we get 
\begin{equation*}
\sfmx[\mP^T (\mW_K^i \mX)^T (\mW_Q^i \mX) \mP] = \mP^T \sfmx[ (\mW_K^i \mX)^T (\mW_Q^i \mX) ] \mP.
\end{equation*}
Then,
\begin{align*}
    \attn(\mX \mP) &= \mX \mP + \sum_{i=1}^h \mW_O^i (\mW_V^i \mX \mP) \cdot \mP^T \sfmx [ (\mW_K^i \mX)^T (\mW_Q^i \mX) ] \mP = \attn(\mX) \mP,
\end{align*}
where we used $\mP \mP^T = \mI$. Permutation equivariance of the token-wise feed-forward layer can be shown similarly:
\begin{align*}
    \ff(\mX \mP) &= \attn(\mX)\mP + \mW_2 \cdot \relu (\mW_1 \cdot \attn(\mX)\mP + \vb_1 \vone_n^T \mP) + \vb_2 \vone_n^T \mP\\
    &= \attn(\mX)\mP + \mW_2 \cdot \relu (\mW_1 \cdot \attn(\mX) + \vb_1 \vone_n^T) \mP + \vb_2 \vone_n^T \mP = \ff(\mX)\mP,
\end{align*}
where $\relu(\mX \mP) = \relu(\mX) \mP$ was used.
This analysis shows that the function class $\mc T^{h,m,r}(\cdot)$ is restricted to permutation equivariant functions.

%%%%%%%%%%%%%%%%%%%%%%%%%%%%%%%%%%%%%%%%%%%%%%%%%%%%%%%

\section{Proof details of Theorem \ref{thm:univapprox}}\label{sec:proof_main}
We first define some additional notation.
For $a, b \in \naturals$ where $a \leq b$, let $[a] = \{1, \dots, a\}$ and $[a:b] = \{a, a+1, \dots, b-1, b\}$.
For $a, b, c \in \reals$ where $b-a > 0$ is an integer multiple of $c > 0$, we write $[a:c:b] \defeq \{a, a+c, a+2c, \dots, b-c, b\}$.

\subsection{Approximating $\mc F_{\rm PE}$ with $\overline{\mc F}_{\rm PE}(\delta)$}
\label{appen:thm2-part1}

\begin{lemma}
\label{lemma:part1}
For any given $f \in \mc F_{\rm PE}$ and $1 \leq p < \infty$, one can find a $\delta^* > 0$ such that $\exists$ $\overline f \in \overline {\mc F}_{\rm PE} (\delta^*)$ which satisfies $\funcdist_p(f, \overline f) \leq \epsilon/3$.
\end{lemma}

\begin{proof}
Since $f: \reals^{d \times n} \to \reals^{d \times n}$ is a continuous function with compact support, the function is uniformly continuous. Since continuity is defined using entry-wise $\ell_p$ norm, and entry-wise $\ell_p$ norm is equivalent to entry-wise $\ell_\infty$ norm when the number of entries are finite, uniform continuity implies that
\begin{equation*}
    \forall \epsilon > 0, \exists \delta > 0 \text { such that } \forall \mX, \mY, \linf{\mX-\mY} < \delta \implies \norm{f(\mX)-f(\mY)}_p < \epsilon.
\end{equation*}
This means that given any $\epsilon/3>0$, we have such a $\delta>0$. Using this $\delta$, we can create a grid $\sG_{\delta}$ and corresponding cubes $\sS_\mL$, as described in the main text.
For any $\mL \in \sG_{\delta}$, we define $\mC_\mL \in \sS_\mL$ to be the center point of the cube $\sS_\mL$. Then, we can define a piece-wise constant approximation $\overline f(\mX) = \sum\nolimits_{\mL \in \sG_{\delta}} f(\mC_\mL) \indic{\mX \in \sS_\mL}$.
Note that, for any $\mX \in \sS_\mL$, we have $\linf{\mX - \mC_\mL} < \delta$, so by uniform continuity, we have $\norm{f(\mX) - \overline f(\mX)}_p = \norm{f(\mX) - f(\mC_\mL)}_p< \epsilon / 3$. 
This proves that $\funcdist_p(f, \overline f) < \epsilon/3$.

As for permutation equivariance, since $f$ is permutation equivariant, we have $f(\mC_\mL \mP) = f(\mC_\mL)\mP$ for any permutation matrix $\mP$. For any $\mX \in \sS_\mL$, we have $\mX \mP \in \sS_{\mL \mP}$, so
\begin{equation*}
    \overline f(\mX \mP)
    = f(\mC_{\mL \mP}) = f(\mC_\mL \mP) = f(\mC_\mL) \mP = \overline f(\mX) \mP.
\end{equation*}
Thus, the approximation $\overline f$ is also permutation equivariant. This proves the lemma.
\end{proof}

\subsection{Approximating $\overline{\mc T}^{2,1,1}$ with $\mc T^{2,1,4}$}
\label{appen:thm2-part3}

\begin{lemma}
\label{lemma:part3}
For each $\overline g \in \overline {\mc T}^{2,1,1}$ and $1 \leq p < \infty$, $\exists$ $g \in \mc T^{2,1,4}$ such that $\funcdist_p(\overline g,g) \leq \epsilon/3$.
\end{lemma}

\begin{proof}
Recall that $T^{h,m,r}$ refers to the class of functions representable with composition of Transformer blocks with $h$ heads of size $m$ in self-attention layers and $r$ hidden nodes in feed-forward layers. The same notation holds for the modified Transformers $\overline {\mc T}^{h,m,r}$.

Note that the softmax operator on a matrix $\mA$ can be made arbitrarily close to hardmax by scaling up $\mA$. That is,
\begin{equation*}
    \sfmx[\lambda \mA] \rightarrow \hdmx[\mA] ~~ \text{ as } \lambda \rightarrow \infty.
\end{equation*}
This means that by scaling up parameters inside $\sfmx$, we can approximate $\hdmx$ arbitrarily closely. Thus, the modified self-attention layers can be approximated with the original self-attention layers of the same number of heads $h$ and head size $m$.

Also, any arbitrary (possibly discontinuous) piecewise linear function $\phi \in \Phi$ can be approximated arbitrarily closely by four $\relu$'s.
Note that $\phi \in \Phi$ as at most three pieces, and at least one of the pieces is constant.
For example, consider the following function $\phi \in \Phi$:
\begin{equation*}
    \phi(t) = 
    \begin{cases}
        b_1 & \text { if } t < c_1,\\
        a_2 t + b_2 & \text { if } c_1 \leq t < c_2,\\
        a_3 t + b_3 & \text { if } c_2 \leq t.
    \end{cases}
\end{equation*}
This function can be approximated by four $\relu$'s, as claimed by the lemma:
\begin{align*}
    \wt \phi(t) =~&
    b_1 + \frac{a_2 c_1 + b_2 - b_1}{\epsilon} \relu(t-c_1+\epsilon)
    + \left ( a_2 - \frac{a_2 c_1 + b_2 - b_1}{\epsilon} \right ) \relu(t-c_1) \\
    &+ \left ( \frac{a_3 c_2 + b_3 - a_2(c_2 - \epsilon) - b_2}{\epsilon} - a_2 \right ) \relu(t - c_2 + \epsilon)\\
    &+ \left ( a_3 - \frac{a_3 c_2 + b_3 - a_2(c_2 - \epsilon) - b_2}{\epsilon} \right ) \relu(t - c_2)\\
    =~&
    \begin{cases}
        b_1 & \text { if } t < c_1-\epsilon,\\
        \frac{a_2 c_1 + b_2 - b_1}{\epsilon} (t - c_1) + a_2c_1 + b_2 & \text { if } c_1-\epsilon \leq t < c_1,\\
        a_2 t + b_2 & \text { if } c_1 \leq t < c_2-\epsilon,\\
        \frac{a_3 c_2 + b_3 - a_2(c_2 - \epsilon) - b_2}{\epsilon} (t - c_2) + a_3 c_2 + b_3 & \text { if } c_2-\epsilon \leq t < c_2,\\
        a_3 t + b_3 & \text { if } c_2 \leq t.
    \end{cases}
\end{align*}
Also, as we make $\epsilon \rightarrow 0$, we can approximate $\phi$ as closely as possible using $\wt \phi$. The cases where the second or third piece is constant can be shown similarly. This means that the modified feed-forward layers (whose activation is $\phi \in \Phi$) with single hidden node can be approximated with the original feed-forward layers ($\relu$) with four hidden nodes.

Thus, given any $\overline{g} \in \overline{\mc T}^{2,1,1}$, there exists a function $g \in {\mc T}^{2,1,4}$ arbitrarily close to $\overline{g}$, by appropriately choosing the parameters to be large enough. This finishes the proof.
\end{proof}

\subsection{Finishing proof of Proposition~\ref{prop:part2}}
%\section{Different roles of attention and feed-forward layers}
\label{sec:appx_proof_prop}

%Recall that our objective is to show that given a piecewise constant function $\overline f \in \overline{\mc F}_{PE}(\delta)$, there exists a modified Transformer network $\overline{g} \in \overline{\mc T}^{2, 1, 1}$ that closely approximates $\overline{f}$. However, as mentioned earlier, the large amount of shared parameters across different tokens in a sequence present in Transformer networks makes this task seemingly difficult. Both the attention and the feed-forward layer weights inside a Transformer block are fixed for all $n$ tokens. Thus, in order to approximate an arbitrary permutation equivariant sequence-to-sequence function, it becomes critical to utilize the contextual information available in the input.

As we have already discussed in Section~\ref{sec:different-roles}, we establish Proposition~\ref{prop:part2} in three steps: 
\begin{enumerate}
\item Given an input $\mX$, a group of feed-forward layers in the modified Transformer network can quantize $\mX$ to an element $\mL$ on the extended grid $\sG^+_{\delta} \defeq \{ -\delta^{-nd}, 0, \delta, \dots, 1-\delta \}^{d \times n}$.
\item Next, a group of self-attention layers in the modified Transformer network can take the input $\mL$ and produce desirable {\em contextual mappings} $q(\mL)$ such that, for $\mL$ and $\tilde{\mL}$, that are not permutation of each other, all the elements in $q(\mL)$ and $q(\tilde{\mL})$ are distinct.
\item Finally, a group of feed-forward layers in the modified Transformer network can map elements of the contextual embedding $q(\mL)$ to the desirable values, i.e., the output of $\overline{f} \in \overline{\mc F}_{\rm PE}$ on the input $\mX$.
\end{enumerate}
These steps are formally stated in Lemmas~\ref{lemma:quantize}, \ref{lemma:contextmap}, and \ref{lemma:memorize} in the main text. We present the proofs of these lemmas in the subsequent sections.

With the results established in these lemmas, we are now equipped with all the tools necessary to complete the proof of Proposition~\ref{prop:part2}. Let us recall the functions $g_{\rm q}, g_{\rm c}$, and $g_{\rm v}$ from Lemma~\ref{lemma:quantize},~\ref{lemma:contextmap}, and~\ref{lemma:memorize}, respectively. We now show that the (modified) Transformer network $\overline{g} = g_{\rm v}\circ g_{\rm c} \circ g_{\rm q}$ approximates the underlying peicewise constant function $\overline{f} \in \overline{\mc F}_{\rm PE}$ over all points in its support except for a set of of measure $O(\delta^d)$.

Consider a point $\mX \in \sS_\mL \subset [0,1]^{d \times n}$, where $\mL \in \wt {\sG}_{\delta}$. By Lemma~\ref{lemma:quantize}, we have that $g_{\rm q}(\mX) = \mL$. Thus, it follows from Lemmas~\ref{lemma:contextmap} and \ref{lemma:memorize} that
$$
g_{\rm v}\circ g_{\rm c} \circ g_{\rm q}(\mX) = g_{\rm v}\circ g_{\rm c}(\mL) = \begin{bmatrix}
    g_{\rm v}^{\rm tkn} (g_{\rm c}(\mL)_{\cdot, 1}) &
    g_{\rm v}^{\rm tkn} (g_{\rm c}(\mL)_{\cdot, 2}) &
    \cdots &
    g_{\rm v}^{\rm tkn} (g_{\rm c}(\mL)_{\cdot, n})
    \end{bmatrix}
    =
    \mA_\mL.
$$
On the other hand, any point $\mX \in \bigcup_{\mL \in \sG_{\delta} \setminus \wt {\sG}_{\delta}} \sS_\mL \cup (\reals^{d\times n} \setminus [0,1]^{d \times n})$ is mapped by $g_{\rm q}$ to $\mL \in \sG^+_{\delta} \setminus \wt {\sG}_{\delta}$; as a result, we get $g_{\rm v}\circ g_{\rm c} \circ g_{\rm q}(\mX) = g_{\rm v} \circ g_{\rm c} (\mL) = \vzero$.

Therefore, we have $\overline{g}(\mX) = g_{\rm v} \circ g_{\rm c} \circ g_{\rm q}(\mX) = \mA_\mL = \overline f(\mX)$ for $\mX \in \bigcup_{\mL \in \wt {\sG}_{\delta}} \sS_\mL$, and $\vzero$ everywhere else. Recall that $\overline{f}$ has its compact support in $[0,1]^{d}$, thus bounded; i.e., there exists $B \geq 0$ such that $\norms{\overline f(\mX)}_p \leq B$. The modified Transformer network $\overline{g}$ takes the same value as $\overline{f}$ on all points in $[0,1]^{d}$ except for a set $\bigcup_{\mL \in \sG_{\delta} \setminus \wt {\sG}_{\delta}} \sS_\mL$ that has measure $O(\delta^d)$. This implies that $\funcdist_p(\overline{f}, \overline{g}) \leq (B^p \delta^d)^{1/p} = O(\delta^{d/p})$.

\subsection{Proof of Lemma~\ref{lemma:quantize}}
The proof strategy is simple; using $\frac{1}{\delta}+1$ token-wise feed-forward layers, we implement the quantization function $g_{\rm q}^{\rm ent}$ that works on the first row of the input. Then stack another $\frac{1}{\delta}+1$ layers that quantizes the second row, and so on.

Given input $\mX$, we first start by clipping $\mX_{1,:}$ in the set $(-\infty, 0) \cup [1, +\infty)$ and mapping the intervals to $-\delta^{-nd}$. This can be done by the following layer:
\begin{equation*}
    \mZ \mapsto \mZ + \ve^{(1)} \phi( (\ve^{(1)})^T \mZ ),
    ~~
    \phi(t) = 
    \begin{cases}
        -t - \delta^{-nd} & \text{ if } t < 0 \text{ or } t \geq 1,\\
        0 & \text{ otherwise.}
    \end{cases}
\end{equation*}

Next, add $1/\delta$ layers of the following form, for $k = 0, \delta, \dots, 1-\delta$.
\begin{equation*}
    \mZ \mapsto \mZ + \ve^{(1)} \phi( (\ve^{(1)})^T \mZ - k\delta \vone_n^T),
    ~~
    \phi(t) = 
    \begin{cases}
    0 & t < 0 \text{ or } t \geq \delta\\
    -t & 0 \leq t < \delta.
    \end{cases}
\end{equation*}
Each layer quantizes $\mX_{1,:}$ in $[k\delta, k\delta+\delta)$ to $k\delta$, without modifying other intervals.

Note that both $\phi$'s used in this construction are piecewise linear functions with three pieces, and at least one of them are constant. Thus, both $\phi$'s are in $\Phi$. We can repeat the same thing for the other rows, and at the end we will get a map from $\reals^{d \times n}$ to $\sG^+_{\delta}$.


\subsection{Proof of Lemma~\ref{lemma:contextmap}}
\label{sec:proof-contextmap}
%Also, for $a, b, c \in \reals$ where $b-a > 0$ is an integer multiple of $c > 0$, we write $[a:c:b] \defeq \{a, a+c, a+2c, \dots, b-c, b\}$.

\paragraph{Selective shift operation.}
Before starting the proof, we first describe the key component of our proof, which we refer to the \emph{selective shift operation}.
Consider the following function, which can be expressed with a multiplicative attention head, with head size $m = 1$ and hardmax $\hdmx$:
\begin{align*}
    \psi(\mZ; b_Q) &= \ve^{(1)} \vu^T \mZ \hdmx [ (\vu^T \mZ)^T (\vu^T \mZ - b_Q \vone_n^T) ]
\end{align*}
where $\vu \in \reals^d$ is a vector that we will choose later, and $\ve^{(1)} = (1, 0, 0, \dots, 0) \in \reals^d$ is the standard basis vector.

To see what this function computes, first consider the $j$-th column of the attention score matrix: $(\vu^T \mZ)^T (\vu^T \mZ_{:,j} - b_Q)$. Note that, if $\vu^T \mZ_{:,j}>b_Q$, $\hdmx$ will calculate $\argmax$ of $\vu^T \mZ$, whereas if $\vu^T \mZ_{:,j}<b_Q$, it will calculate $\argmin$. Therefore, the $(1,j)$-th entry of $\psi(\mZ; b_Q) \in \reals^{d \times n}$ can be written as
\begin{equation*}
\psi(\mZ; b_Q)_{1, j} = 
\vu^T \mZ \hdmx [ (\vu^T \mZ)^T (\vu^T \mZ_{:,j} - b_Q) ] =
\begin{cases}
    \max_k \vu^T \mZ_{:,k} & \text{ if } \vu^T \mZ_{:,j} > b_Q, \\
    \min_k \vu^T \mZ_{:,k} & \text{ if } \vu^T \mZ_{:,j} < b_Q,
\end{cases}
\end{equation*}
for $j \in [n]$. Note that due to $\ve^{(1)}$, all rows of $\psi(\mZ; b_Q)$ except the first row are zero.
From this observation, one can define a function parametrized by $b_Q$ and $b'_Q$, where $b_Q < b'_Q$, which consists of two attention heads:
\begin{align*}
    &\Psi(\mZ; b_Q, b'_Q) \defeq \psi(\mZ; b_Q) - \psi(\mZ; b'_Q),\\
    &\Psi(\mZ; b_Q, b'_Q)_{1, j} =
    \begin{cases}
        \max_k \vu^T \mZ_{:,k} - \min_k \vu^T \mZ_{:,k} & \text{ if } b_Q < \vu^T \mZ_{:,j} < b'_Q, \\
        0 & \text{ if } \vu^T \mZ_{:,j} < b_Q \text{ or } \vu^T \mZ_{:,j} > b'_Q.
    \end{cases}
\end{align*}
What this means is that, if we define an attention layer of the form $\mZ \mapsto \mZ + \Psi(\mZ;b_Q,b'_Q)$, then any column $\mZ_{:,j}$ satisfying $\vu^T \mZ_{:,j} \in (b_Q,b'_Q)$ is shifted up in its first coordinate $\mZ_{1,j}$ by $\max_k \vu^T \mZ_{:,k} - \min_k \vu^T \mZ_{:,k}$, \textbf{while all the other coordinates stay untouched}. We call this the selective shift operation, because we can choose $b_Q$ and $b'_Q$ to selectively shift certain entries of the input.

\paragraph{Bijective column id mapping.}
Recall that the input to this step is from the range of $g_{\rm q}$ (Lemma~\ref{lemma:quantize}), which is $\sG^+_{\delta} = \{ -\delta^{-nd}, 0, \delta, \dots, 1-\delta \}^{d \times n}$. Now consider $\mL \in \sG^+_{\delta}$ and $\vu = (1, \delta^{-1}, \delta^{-2}, \dots, \delta^{-d+1})$.

For any $j \in [n]$, it is easy to check two following facts:
\begin{enumerate}
\item If $\mL_{i,j} \neq -\delta^{-nd}$ for all $i \in [d]$, i.e., $\mL_{:,j} \in \{0, \delta, \dots, 1-\delta\}^d$, then $\vu^T \mL_{:,j} \in [0:\delta:\delta^{-d+1} - \delta]$, and the map $\mL_{:,j} \mapsto \vu^T \mL_{:,j}$ from $\{0, \delta, \dots, 1-\delta\}^d$ to $[0:\delta:\delta^{-d+1} - \delta]$ is a bijection.
\item If there exists $i \in [d]$ such that $\mL_{i,j} = -\delta^{-nd}$, then $\vu^T \mL_{:,j} \leq -\delta^{-nd} + \delta^{-d+1} - 1 < 0$.
\end{enumerate}
Therefore, one can say that $\vu^T \mL_{:,j}$ gives the ``column id'' for each possible value of $\mL_{:,j} \in \{0, \delta, \dots, 1-\delta\}^d$.

The rough idea of the construction is to apply the selective shift operation to each column id, by setting $\vu$ in the definition of $\Psi(\cdot)$ to be $(1, \delta^{-1}, \delta^{-2}, \dots, \delta^{-d+1})$ and choosing $b_Q = l - \delta/2$ and $b'_Q = l + \delta/2$ for each $l \in [0:\delta:\delta^{-d+1} - \delta]$.
More concretely, we stack $(1/\delta)^d$ attention layers, with attention parts $\delta^{-d} \Psi(\cdot;l-\delta/2,l+\delta/2)$ for each $l \in [0:\delta:\delta^{-d+1} - \delta]$, in increasing order of $l$.
After that, we add an extra single-head attention layer with attention part $\delta^{-(n+1)d} \psi(\cdot;0)$.

We now divide possible input values $\mL \in \sG^+_{\delta}$ into three disjoint categories, and show how these layers change the input values at the end of all the layers. Recall the hierarchy $\wt {\sG}_{\delta} \subset \sG_{\delta} \subset \sG^+_{\delta}$. The categories are defined as follows: 
\begin{enumerate}
    \item $\mL \in \wt {\sG}_{\delta}$. All entries are between $0$ and $1-\delta$, and all columns are unique.
    \item $\mL \in \sG_{\delta} \setminus \wt {\sG}_{\delta}$. All entries are between $0$ and $1-\delta$, but there are duplicate columns.
    \item $\mL \in \sG^+_{\delta} \setminus \sG_{\delta}$. The point has at least one entry that equals to $-\delta^{-nd}$.
\end{enumerate}

\subsubsection{Category 1}
\label{sec:proof-category1}
In Category~1, we have $\mL \in \wt {\sG}_{\delta}$.
Let $l_j \defeq \vu^T \mL_{:,j}$.
Due to permutation equivariance, we can assume without loss of generality that $l_j$'s are in increasing order: $l_1 < l_2 < \cdots < l_n$.
The first $(1/\delta)^d$ layers sweep the set $[0:\delta:\delta^{-d+1}-\delta]$ and apply selective shift operation on each element in the set. This means that selective shift operation will be applied to $l_1$ first, then $l_2$, and then $l_3$, and so on, regardless of the specific values of $l_j$'s.

\paragraph{First shift operation.}
In the first selective shift operation, the $(1,1)$-th entry of $\mL$ ($\emL_{1,1}$) is shifted by the operation, while the other entries are left untouched. The updated value $\wt \emL_{1,1}$ is
\begin{align*}
    \wt \emL_{1,1} 
    = \emL_{1,1} + \delta^{-d} (\max\nolimits_k \vu^T \mL_{:,k} - \min\nolimits_k \vu^T \mL_{:,k}) 
    = \emL_{1,1} + \delta^{-d} (l_n - l_1).
\end{align*}
Therefore, after the operation, the output of the layer is $\begin{bmatrix} \wt \mL_{:,1} & \mL_{:,2} & \dots & \mL_{:,n} \end{bmatrix}$, and the new value of the first column $\wt \mL_{:,1}$ results in
\begin{align*}
    \vu^T \wt \mL_{:,1}
    = \wt \emL_{1,1} + \sum_{i=2}^d \delta^{-i+1} \emL_{i,1}
    = \emL_{1,1} + \delta^{-d} (l_n - l_1) + \sum_{i=2}^d \delta^{-i+1} \emL_{i,1}
    = l_1 + \delta^{-d} (l_n - l_1).
\end{align*}
Let us denote the updated ``column id'' $\vu^T \wt \mL_{:,1}$ as $\wt l_1$.
% Since $l_j \in [0:\delta:\delta^{-d+1}-\delta]$ we have
% \begin{align*}
%     \delta^{-d+1} = 0 + \delta^{-d} \cdot \delta \leq \wt l_1 \defeq l_1 + \delta^{-d} (l_n - l_1) \leq \delta^{-d+1} - \delta + \delta^{-d}(\delta^{-d+1}-\delta) = \delta^{-2d+1}-\delta,
% \end{align*}
% so in particular, $l_n < \wt l_1$. 
We can show that $l_n < \wt l_1$, because
\begin{align*}
    \wt l_1 \defeq l_1 + \delta^{-d} (l_n - l_1)
    \geq 0 + \delta^{-d} \cdot \delta
    = \delta^{-d + 1} > l_n.
\end{align*}
Therefore, after updating, 
\begin{equation*}
    \max \vu^T \begin{bmatrix} \wt \mL_{:,1} & \mL_{:,2} & \dots & \mL_{:,n} \end{bmatrix}
    = \max \{ \wt l_1, l_2, \dots, l_n \} = \wt l_1,
\end{equation*}
and the new minimum is $l_2$.

\paragraph{Second shift operation.}
The second selective shift operation is applied to $l_2$, by which only one entry $\emL_{1,2}$ will be shifted. The updated value $\wt \emL_{1,2}$ is
\begin{equation*}
    \wt \emL_{1,2}
    = \emL_{1,2} + \delta^{-d} (\wt l_1 - l_2)
    = \emL_{1,2} + \delta^{-d} (l_1 - l_2) + \delta^{-2d} (l_n - l_1).
\end{equation*}
After updating, the new inner product of $\vu$ and $\wt \mL_{:,2}$ results in
\begin{equation*}
    \wt l_2 \defeq \vu^T \wt \mL_{:,2} = l_2 + \delta^{-d} (l_1 - l_2) + \delta^{-2d} (l_n - l_1).% = (1-\delta^{-d}) l_2 + (\delta^{-d} - \delta^{-2d}) l_1 + \delta^{-2d} l_n.
\end{equation*}
% In a similar way as $\wt l_1$, we can show the following bounds:
% \begin{align*}
%     \wt l_2 &\defeq (1-\delta^{-d}) l_2 + (\delta^{-d} - \delta^{-2d}) l_1 + \delta^{-2d} l_n 
%     \geq (1-\delta^{-d}) l_2 + (\delta^{-d} - \delta^{-2d}) l_2 + \delta^{-2d} l_n \\
%     &= (1-\delta^{-2d}) l_2 + \delta^{-2d} l_n 
%     = l_2 + \delta^{-2d}(l_n - l_2) \geq 0 + \delta^{-2d}\cdot \delta = \delta^{-2d+1},\\
%     \wt l_2 &\defeq (1-\delta^{-d}) l_2 + (\delta^{-d} - \delta^{-2d}) l_1 + \delta^{-2d} l_n 
%     \leq (1-\delta^{-d}) l_1 + (\delta^{-d} - \delta^{-2d}) l_1 + \delta^{-2d} l_n\\
%     &= l_1 + \delta^{-2d}(l_n - l_1)
%     \leq \delta^{-d+1} - \delta + \delta^{-2d} (\delta^{-d+1} - \delta) 
%     \leq \delta^{-3d+1} - \delta.
% \end{align*}
We can show that $\wt l_1 < \wt l_2$, because
\begin{align*}
    &l_1 + \delta^{-d} (l_n - l_1) < l_2 + \delta^{-d} (l_1 - l_2) + \delta^{-2d} (l_n - l_1)\\
    \iff~&
    (\delta^{-d}-1) (l_2 - l_1) < \delta^{-d}(\delta^{-d}-1) (l_n - l_1),
\end{align*}
and the last inequality is true because $\delta^{-d} > 1$ and $l_n > l_2$.
Since we have $\wt l_1 < \wt l_2$, and the new maximum in $\vu^T \begin{bmatrix} \wt \mL_{:,1} & \wt \mL_{:,2} & \mL_{:,3} & \dots & \mL_{:,n} \end{bmatrix}$ is now $\wt l_2$, and the new minimum is $l_3$.

\paragraph{Repeating the process.}
More generally, we can repeat this process, and show that the $j$-th shift operation shifts $\emL_{1,j}$ by $\delta^{-d} (\wt l_{j-1} - l_j)$, and results in the new column id
\begin{equation*}
    \wt l_j \defeq 
    \vu^T \wt \mL_{:,j}
    =
    l_j + \sum_{k=1}^{j-1} \delta^{-kd}(l_{j-k}-l_{j-k+1}) + \delta^{-jd} (l_n-l_1).
\end{equation*}
In the general case, $\wt l_{j-1} < \wt l_j$ holds $j = [2:n]$, because
\begin{align*}
     &\wt l_{j-1} = l_{j-1} + \sum_{k=2}^{j-1} \delta^{-kd+d}(l_{j-k}-l_{j-k+1}) + \delta^{-(j-1)d} (l_n-l_1)\\
     &< \wt l_j = l_j + \sum_{k=1}^{j-1} \delta^{-kd}(l_{j-k}-l_{j-k+1}) + \delta^{-jd} (l_n-l_1)\\
     \iff~&
     \sum_{k=1}^{j-1} \delta^{-kd+d} (\delta^{-d}-1)(l_{j-k+1}-l_{j-k}) <
     \delta^{-(j-1)d}(\delta^{-d}-1)(l_n-l_1),
\end{align*}
and the last inequality holds because
\begin{equation*}
    \delta^{-(j-1)d}(l_n-l_1) > \delta^{-(j-1)d} \sum_{k=1}^{j-1}(l_{j-k+1} - l_{j-k}) > \sum_{k=1}^{j-1} \delta^{-kd+d} (l_{j-k+1}-l_{j-k}).
\end{equation*}
Therefore, after the $j$-th selective shift operation, $\wt l_j$ is the new maximum among $\{\wt l_1, \dots, \wt l_j, l_{j+1}, \dots, l_n\}$ and $l_{j+1}$ is the new minimum, which makes us possible to continue the process until the $n$-th operation.
% \begin{equation*}
% %\label{eq:tildelk}
%     \wt l_j \defeq \sum_{k=0}^{j-1} \delta^{-kd}(1-\delta^{-d}) l_{j-k} + \delta^{-jd} l_n.
% \end{equation*}
% For $1 \leq j < n$, $\wt l_j$ satisfies the following bounds
% \begin{align*}
%     \wt l_j
%     &\geq \sum_{k=0}^{j-1} \delta^{-kd}(1-\delta^{-d}) l_{j} + \delta^{-jd} l_n
%     = (1-\delta^{-jd})l_{j} + \delta^{-jd} l_n \\
%     &= l_{j} + \delta^{-jd} (l_n - l_j)
%     \geq 0 + \delta^{-jd} \cdot \delta = \delta^{-jd + 1},\\
%     \wt l_j 
%     &\leq \sum_{k=0}^{j-1} \delta^{-kd}(1-\delta^{-d}) l_{1} + \delta^{-jd} l_n
%     = (1-\delta^{-jd})l_{1} + \delta^{-jd} l_n \\
%     &= l_1 + \delta^{-jd}(l_n - l_1) 
%     \leq \delta^{-d+1} - \delta + \delta^{-jd} (\delta^{-d+1} - \delta)
%     \leq \delta^{-(j+1)d+1} - \delta.
% \end{align*}

\paragraph{After $n$ shift operations.}
As a result, after the whole sweep from $0$ to $\delta^{-d+1}-\delta$ by the first $(1/\delta)^d$ layers, a total of $n$ shift operations are applied, and the input $\mL$ is mapped to a new point $\wt \mL$, where $\vu^T \wt \mL = \begin{bmatrix}\wt l_1 & \wt l_2 & \dots & \wt l_n \end{bmatrix}$ and $\wt l_1 < \wt l_2 < \dots < \wt l_n$.

We can now prove the following technical lemma, whose proof is deferred to Appendix~\ref{sec:proof-sublemma:1}:
\begin{lemma}
\label{sublemma:1}
After $n$ shift operations, $\wt l_n = \vu^T \wt \mL_{:,n}$ satisfies the following bounds:
\begin{equation*}
    \delta^{-(n-1)d+1}(\delta^{-d} - 1) 
    \leq \wt l_n 
    \leq \delta^{-nd + 1}(\delta^{-d}-1) - \delta (\delta^{-d} - 1)^2.
\end{equation*}
Also, the map from $\begin{bmatrix} l_1 & l_2 & \cdots & l_n \end{bmatrix} \in [0:\delta:\delta^{-d+1}-\delta]$ (where $l_1 < l_2 < \dots < l_n$) to $\wt l_n$ is one-to-one.
\end{lemma}

\paragraph{Global shifting by the last layer.}
As mentioned earlier, after this sweep, there is another attention layer with attention part $\delta^{-(n+1)d} \psi(\cdot;0)$.
Since $0<\wt l_1 < \cdots < \wt l_n$, what it does to $\wt \mL$ is that it adds $\delta^{-(n+1)d} \max_k \vu^T \wt \mL_{:,k} = \delta^{-(n+1)d} \wt l_n$ to each entry in the first row of $\wt \mL$. The output of this layer is defined to be the function $g_{\rm c}(\mL)$.

Now, in summary, for any $\mL \in \wt {\sG}_{\delta}$, $i \in [d]$, and $j \in [n]$, we have
\begin{equation*}
    g_{\rm c}(\mL)_{i,j} = 
    \begin{cases}
        \emL_{1,j} + \sum_{k=1}^{j-1} \delta^{-kd}(l_{j-k}-l_{j-k+1}) + \delta^{-jd} (l_n-l_1) + \delta^{-(n+1)d} \wt l_n & \text { if } i = 1,\\
        \emL_{i,j} & \text { if } i \in [2,d],
    \end{cases}
\end{equation*}
and for any $\mL \in \wt {\sG}_{\delta}$ and $j \in [n]$,
\begin{equation*}
    \vu^T g_{\rm c}(\mL)_{:,j} = \wt l_{j} + \delta^{-(n+1)d} \wt l_n.
\end{equation*}

\paragraph{Checking Properties~\ref{lemma:contextmap}.\ref{cond:1} and  \ref{lemma:contextmap}.\ref{cond:2}.}
Given this result so far, it is now left to check if the constructed network is really a permutation equivariant contextual mapping, i.e., if it satisfies Properties~\ref{lemma:contextmap}.\ref{cond:1} and \ref{lemma:contextmap}.\ref{cond:2} in Lemma~\ref{lemma:contextmap}.

First, for any $\mL \in \wt {\sG}_{\delta}$, Property~\ref{lemma:contextmap}.\ref{cond:1} holds because we already know $\wt l_1 < \wt l_2 < \dots < \wt l_n$, so they are all distinct.
As for Property~\ref{lemma:contextmap}.\ref{cond:2}, note that the upper bound on $\wt l_n$ from Lemma~\ref{sublemma:1} also holds for other $\wt l_j$'s, so
\begin{equation*}
\vu^T g_{\rm c}(\mL)_{:,j} \in [\delta^{-(n+1)d} \wt l_n, \delta^{-(n+1)d} \wt l_n + \delta^{-(n+1)d+1}),
\end{equation*}
for all $j \in [n]$. Now, from Lemma~\ref{sublemma:1}, two $\mL, \mL' \in \wt {\sG}_{\delta}$ (that are not permutations of each other) map to different $\wt l_n$ and $\wt l'_n$, and they differ at least by $\delta$. This means that two intervals $[\delta^{-(n+1)d} \wt l_n, \delta^{-(n+1)d} \wt l_n + \delta^{-(n+1)d+1})$ and $[\delta^{-(n+1)d} \wt l'_n, \delta^{-(n+1)d} \wt l'_n + \delta^{-(n+1)d+1})$ are guaranteed to be disjoint, so the entries of $\vu^T g_{\rm c}(\mL)$ and $\vu^T g_{\rm c}(\mL')$ are all distinct. This proves Property~\ref{lemma:contextmap}.\ref{cond:2}.

Therefore, we finished showing that the map $g_{\rm c}(\cdot)$ we constructed using $(1/\delta)^d+1$ attention layers implements a permutation equivariant contextual mapping on $\wt {\sG}_{\delta}$.

\paragraph{Checking Property~\ref{lemma:contextmap}.\ref{cond:3}.}
It is now left to check if the map $g_{\rm c}$ satisfies the other properties.
At this point, we can check Property~\ref{lemma:contextmap}.\ref{cond:3}. 
From 
$\vu^T g_{\rm c}(\mL)_{:,j} \in [\delta^{-(n+1)d} \wt l_n, \delta^{-(n+1)d} \wt l_n + \delta^{-(n+1)d+1})$ and Lemma~\ref{sublemma:1}, we can show that for any $\mL \in \wt {\sG}_{\delta}$, we have
\begin{align*}
    \delta^{-2nd+1}(\delta^{-d} - 1) 
    \leq \vu^T g_{\rm c}(\mL)_{:,j} 
    &< \delta^{-(n+1)d}(\delta^{-nd + 1}(\delta^{-d}-1) - \delta (\delta^{-d} - 1)^2) + \delta^{-(n+1)d+1}\\
    &\leq \delta^{-(2n+1)d+1}(\delta^{-d}-1),
\end{align*}
where we used $\delta^{-1} \geq 2$.
This proves that all $\vu^T g_{\rm c}(\mL)_{:,j}$ are between $t_l = \delta^{-2nd+1}(\delta^{-d} - 1)$ and $t_r = \delta^{-(2n+1)d+1}(\delta^{-d}-1)$.
For the remaining input points $\mL \in \sG^+_{\delta} \setminus \wt {\sG}_{\delta}$, we will check that $\vu^T g_{\rm c}(\mL)_{:,j}$ is outside the interval $[t_l , t_r]$ (Property~\ref{lemma:contextmap}.\ref{cond:4}).

\subsubsection{Category 2}
In Category~2, we have $\mL \in \sG_{\delta}\setminus \wt {\sG}_{\delta}$. Here, all entries are between $0$ and $1-\delta$, but there are duplicate columns.
Again, let $l_j \defeq \vu^T \mL_{:,j}$, and assume without loss of generality that $l_1 \leq l_2 \leq \dots \leq l_n$. 
For the input $\mL$ in Category~2, there exist some $j, j' \in [n]$, $j \neq j'$, such that $l_j = l_{j'}$.
This means that when the input passes through the attention layer $\delta^{-d}\Psi(\cdot; l_j - \delta/2, l_j + \delta/2)$, the selective shift operation for $l_j$ is applied to both $j$-th and $j'$-th columns; the two columns are coupled together.
More generally, suppose we have $n' < n$ distinct columns.

\paragraph{If $n'= 1$.}
In the extreme case of $n' = 1$, we have $\max_j l_j = \min_j l_j$, so the selective shift operation applied at $l_j$ does not shift the entry at all; therefore, at the end of the first $(1/\delta)^d$ attention layers, $\wt \mL = \mL$.

\paragraph{If $1 < n' \leq n - 1$.}
When $1 < n' \leq n-1$, let the $n'$ distinct values of $l_j$'s be $l'_1, \dots, l'_{n'}$.
The shift operation is applied $n'$ times, to $l'_1, \dots, l'_{n'}$, and shifts one or more entries at a time. After the first $(1/\delta)^d$ layers, the output $\wt \mL$ has $n'$ distinct $\wt l_j = \vu^T \wt \mL_{:,j}$, $0 \leq \wt l_1 \leq \wt l_2 \leq \dots \leq \wt l_n$, whose distinct values are the same as the numbers we get when we apply shift operations to a length-$n'$ sequence $\begin{bmatrix} l'_1 & \dots & l'_{n'} \end{bmatrix}$.
Then, applying the same calculations from Category~1 shows that
\begin{equation*}
    \wt l_n = 
    \vu^T \wt \mL_{:,n}
    =
    l'_{n'} + \sum_{k=1}^{n'-1} \delta^{-kd}(l'_{n'-k}-l'_{n'-k+1}) + \delta^{-n'd} (l'_{n'}-l'_1),
\end{equation*}
and it follows from the upper bound in Lemma~\ref{sublemma:1} that
\begin{equation*}
    \wt l_n
    \leq \delta^{-n'd + 1}(\delta^{-d}-1) - \delta (\delta^{-d} - 1)^2
    < \delta^{-(n-1)d + 1}(\delta^{-d}-1).
\end{equation*}
Note that the RHS matches the lower bound in Lemma~\ref{sublemma:1}.
This implies that the value of $\wt l_n$ calculated from the input $\mL \in \sG_{\delta}\setminus \wt {\sG}_{\delta}$ (Category~2) is always strictly less (by at least $\delta$) than that calculated from $\mL \in \wt {\sG}_{\delta}$ (Category~1).

\paragraph{Checking Property~\ref{lemma:contextmap}.\ref{cond:4}.}
After the global shifting by the last layer with attention part $\delta^{-(n+1)d} \psi(\cdot;0)$, we get the output $g_{\rm c}(\mL)$ which satisfies
\begin{align*}
    \vu^T g_{\rm c}(\mL)_{:,j} = \wt l_{j} + \delta^{-(n+1)d} \wt l_n
    &\leq (\delta^{-(n+1)d} + 1) (\delta^{-(n-1)d + 1}(\delta^{-d}-1) - \delta (\delta^{-d} - 1)^2)\\
    &< \delta^{-2nd+1}(\delta^{-d}-1) \eqdef t_l.
\end{align*}
where the RHS is a lower bound on possible values of $\vu^T g_{\rm c}(\mL)_{:,j}$ for $\mL \in \wt {\sG}_{\delta}$ (Category~1).
This means that the entries of $\vu^T g_{\rm c}(\mL)$ for Category~2 are outside $[t_l, t_r]$, which satisfies Property~\ref{lemma:contextmap}.\ref{cond:4}.

\subsubsection{Category~3}
In Category~3, we have $\mL \in \sG^+_{\delta} \setminus \sG_{\delta}$; the point $\mL$ has at least one entry that equals to $-\delta^{-nd}$.
Let $l_j \defeq \vu^T \mL_{:,j}$, and recall that whenever a column $\mL_{:,j}$ has an entry that equals to $-\delta^{-nd}$, we have $l_j = \vu^T \mL_{:,j} \leq -\delta^{-nd} + \delta^{-d+1} - 1 < 0$.
Assume without loss of generality that $l_1 \leq l_2 \leq \dots \leq l_n$.

Recall that the selective shift operation is applied to each element of $[0:\delta:\delta^{-d+1}-\delta]$, not to negative values.
In case of Category~3, we have $\min_k \vu^T \mL_{:,k} = l_1 < 0$, and $l_1$ never gets shifted upwards, so it remains as the minimum for the whole time.

\paragraph{If all $l_j$'s are negative.}
In case where all $l_j$'s are negative, selective shift operation never changes the input $\mL$, so we get $\wt \mL = \mL$. Since we have $\vu^T \wt \mL < \vzero_n^T$ (entry-wise), the last layer with attention part $\delta^{-(n+1)d} \psi(\cdot;0)$ adds $\delta^{-(n+1)d} \min_k \vu^T \wt \mL_{:,k} < 0$ to each entry in the first row of $\wt \mL$, further pushing it to the negative side. Therefore, the final output $g_{\rm c}(\mL)$ satisfies $\vu^T g_{\rm c}(\mL) < \vzero_n^T < t_l\vone_n^T$.

\paragraph{If not all $l_j$'s are negative.}
Now consider the case where at least one $l_j$ is positive. Let $i$ be the index that satisfies $l_{i-1} < 0 \leq l_{i}$. Then, selective shift operation does not affect $l_1, \dots, l_{i-1}$, and then it shifts $l_{i}$ by
\begin{equation*}
    \delta^{-d}(\max_k \vu^T \mL_{:,k} - \min_k \vu^T \mL_{:,k})
    = \delta^{-d} (l_n - l_1)
    \geq \delta^{-d} (0 + \delta^{-nd} - \delta^{-d+1} + 1)
    \geq \delta^{-(n+1)d+1},
\end{equation*}
where we used $\delta^{-1} \geq 2$ at the last inequality.
The next shift operations shift $l_{i+1}, \dots, l_n$ by even larger amount, so at the end of the first $(1/\delta)^d$ layers, we have 
$\delta^{-(n+1)d+1} \leq \wt l_{i} \leq \dots \leq \wt l_n$,
while $\wt l_j = l_j < 0$ for $j \in [i-1]$.

\paragraph{Shifts by the last layer.}
Here, the last layer with attention part $\delta^{-(n+1)d} \psi(\cdot;0)$ acts differently for negative and positive $\wt l_j$'s. 
For negative $\wt l_j$'s, it adds $\delta^{-(n+1)d} \min_k \wt l_{k} = \delta^{-(n+1)d} l_1 < 0$ to $\wt l_1, \dots, \wt l_{i-1}$, pushing them further to the negative side.
For positive $\wt l_j$'s, the layer adds $\delta^{-(n+1)d} \max_k \wt l_{k} = \delta^{-(n+1)d} \wt l_n \geq \delta^{-(2n+2)d+1}$ to $\wt l_{i}, \dots, \wt l_{n}$, so that they are all greater than or equal to $\delta^{-(2n+2)d+1}$. Note that $\delta^{-(2n+2)d+1} > t_r$.

\paragraph{Checking Property~\ref{lemma:contextmap}.\ref{cond:4}.}
Therefore, in both cases, we can see that the final output $g_{\rm c}(\mL)$ satisfies $\vu^T g_{\rm c}(\mL)_{:,j} \notin [t_l, t_r]$, for all $j \in [n]$.
This completes the verification of Property~\ref{lemma:contextmap}.\ref{cond:4}.

\subsubsection{Proof of Lemma~\ref{sublemma:1}}
\label{sec:proof-sublemma:1}
Proof of lower and upper bounds on $\wt l_n$ are straightforward:
\begin{align*}
    \wt l_n
    &\defeq l_n + \sum_{k=1}^{n-1} \delta^{-kd}(l_{n-k}-l_{n-k+1}) + \delta^{-nd} (l_n-l_1)\\
    &\geq \delta^{-(n-1)d}\sum_{k=1}^{n-1} (l_{n-k}-l_{n-k+1}) + \delta^{-nd} (l_n-l_1)
    = (\delta^{-nd}-\delta^{-(n-1)d}) (l_n-l_1)\\
    &\geq \delta^{-(n-1)d+1}(\delta^{-d} - 1),\\
    \wt l_n
    &\leq l_n + \delta^{-d} (l_1 - l_n) + \delta^{-nd}(l_n - l_1)
    \leq \delta^{-d+1} - \delta + (\delta^{-nd} - \delta^{-d}) (\delta^{-d+1} - \delta)\\
    &=\delta^{-nd + 1}(\delta^{-d}-1) - \delta (\delta^{-2d} - 2\delta^{-d} + 1)
    = \delta^{-nd + 1}(\delta^{-d}-1) - \delta (\delta^{-d} - 1)^2.
\end{align*}

For one-to-one property of the map, consider $\begin{bmatrix} l_1 & l_2 & \cdots & l_n \end{bmatrix}$ and $\begin{bmatrix} l'_1 & l'_2 & \cdots & l'_n \end{bmatrix}$ with increasing entries, which are mapped to $\wt l_n$ and $\wt l'_n$, respectively.
Suppose $\wt l_n = \wt l'_n$. By definition,
\begin{align*}
    \wt l_n - \wt l'_n
    =& (l_n - l'_n) + \delta^{-d} (l_{n-1} - l_n - l'_{n-1} + l'_n) 
    + \delta^{-2d} (l_{n-2} - l_{n-1} - l'_{n-2} + l'_{n-1}) + \dots
    \\&+ \delta^{-(n-1)d}(l_1-l_2-l'_1+l'_2) + \delta^{-nd} (l_n-l_1-l'_n+l'_1) = 0.
\end{align*}
Now assume for contradiction that $l_n \neq l'_n$. Then, we have $-\delta^{-d+1}+\delta \leq l_n - l'_n \leq \delta^{-d+1}-\delta$. However, the remaining terms have ``coarse resolution'', and they can never cancel $l_n - l'_n$ and make the sum zero, because for example, $\delta^{-d} (l_{n-1} - l_n - l'_{n-1} + l'_n)$ can only have values $0, \delta^{-d+1}, -\delta^{-d+1}, 2\delta^{-d+1}, -2\delta^{-d+1}, \dots$. Thus, $l_n = l'_n$ must hold and the first term must be zero.

Similarly, assume that $l_{n-1} \neq l'_{n-1}$. Then, the second term is in the interval $[-\delta^{-2d+1}+\delta^{-d+1}, \delta^{-2d+1}-\delta^{-d+1}]$.
Again, the remaining terms cannot cancel the second term, hence $l_{n-1} = l'_{n-1}$ must hold.
We can proceed this way, and show that $l_j = l'_j$ must hold for all $j \in [n]$, hence proving that the map is one-to-one.


%%%%%%%%%%%%%%%%%%%%%%%%%%%%%%%%%%%%%%%%%%%%%%%%

\subsection{Proof of Lemma~\ref{lemma:memorize}}
Note that $|\sG^+_{\delta}| = (\frac{1}{\delta}+1)^{dn}$, so the image of $g_{\rm c}(\sG^+_{\delta})$ (from Lemma~\ref{lemma:contextmap}) has finite number of distinct real numbers. Let $M$ be the maximum over all these numbers. By construction of $g_{\rm c}$, we know that $M>0$.

To construct a function $g_{\rm v}^{\rm tkn}$ that satisfies the statement of the lemma, we first implement the second part: $g_{\rm v}^{\rm tkn}(g_{\rm c}(\mL)_{:,j}) = \vzero_d$ if $\mL \in \sG^+_{\delta} \setminus \wt \sG_{\delta}$.
Note from Lemma~\ref{lemma:contextmap} that, for any $\mL \in \wt \sG_{\delta}$, we have $\vu^T g_{\rm c}(\mL)_{:,j} \in [t_l, t_r]$ for all $j$, and for any $\mL \in \sG^+_{\delta} \setminus \wt \sG_{\delta}$, $\vu^T g_{\rm c}(\mL)_{:,j} \notin [t_l, t_r]$ for all $j$.
Using this, we add the following feed-forward layer:
\begin{equation*}
    \mZ \mapsto \mZ - (M+1) \vone_n \phi (\vu^T \mZ),
    ~~
    \phi(t)
    =
    \begin{cases}
        0 & \text{ if } t \in [t_l, t_r]\\
        1 & \text{ if } t \notin [t_l, t_r].
    \end{cases}
\end{equation*}
Input to this layer is $g_{\rm c}(\mL)$. If $\mL \in \wt \sG_{\delta}$, then $\phi(\vu^T g_{\rm c}(\mL)) = \vzero_n^T$, so the output stays the same as the input.
If $\mL \in \sG^+_{\delta} \setminus \wt \sG_{\delta}$, then $\phi(\vu^T g_{\rm c}(\mL)) = \vone_n^T$, so all the entries of the input are shifted by $-M-1$, and become strictly negative.

Recall that by definition of $\wt \sG_{\delta}$, all the entries of $g_{\rm c}(\mL)$ for $\mL \in \wt \sG_{\delta}$ are nonnegative. So the next thing to do is mapping all strictly negative entries to zero.
This can be done in a similar way as Lemma~\ref{lemma:quantize}.
For $i \in [d]$, add the following layer:
\begin{equation*}
    \mZ \mapsto \mZ + \ve^{(i)} \phi ((\ve^{(i)})^T \mZ),
    ~~
    \phi(t)
    =
    \begin{cases}
        -t & \text{ if } t < 0\\
        0 & \text{ if } t \geq 0.
    \end{cases}
\end{equation*}
After these $d$ layers, the output for $\mL \in \sG^+_{\delta} \setminus \wt \sG_{\delta}$ is a zero matrix, while the output for $\mL \in \wt \sG_{\delta}$ is $g_{\rm c}(\mL)$.

Now, it is left to map $g_{\rm c}(\mL)$ to $\mA_\mL$, for $\mL \in \wt \sG_{\delta}$. Up to permutation equivariance, each different context $\mL$ maps to $n$ unique numbers $\vu^T g_{\rm c}(\mL)$, which are at least $\delta$ apart from each other. The idea of value mapping is to map each unique number to the corresponding output column.

More precisely, choose any $\overline \mL \in \wt \sG_{\delta}$.
For each value of $\vu^T g_{\rm c}(\overline \mL)_{:,j}$, $j \in [n]$, we add one feed-forward layer
\begin{equation*}
    \mZ \mapsto \mZ + ((\mA_{\overline \mL})_{:,j} - g_{\rm c}({\overline \mL})_{:,j}) \phi (\vu^T \mZ - \vu^T g_{\rm c}({\overline \mL})_{:,j} \vone_n^T),
    ~~
    \phi(t) = 
    \begin{cases}
    0 & t < -\delta/2 \text{ or } t \geq \delta/2,\\
    1 & -\delta/2 \leq t < \delta/2.
    \end{cases}
\end{equation*}
If the input $\mZ$ is a zero matrix, which is the case for $\mL \in \sG^+_{\delta} \setminus \wt \sG_{\delta}$, $\vu^T \mZ = \vzero_n^T$. Since $t_l$ is much larger than $0$, activation is all zero. Thus, zero input matrix remains the same at the output.

If the input $\mZ$ is $g_{\rm c}(\mL)$, where $\mL \in \wt \sG_{\delta}$ is not a permutation of $\overline \mL$, then 
\begin{equation*}
    \phi(\vu^T g_{\rm c}(\mL) - \vu^T g_{\rm c}({\overline \mL})_{:,j} \vone_n^T) = \vzero_n^T,   
\end{equation*}
so $g_{\rm c}(\mL)$ is left untouched.

If some other $\mL$ is a permutation of $\overline \mL$, and $\mL_{:,i} = \overline \mL_{:,j}$, then 
\begin{equation*}
    \phi(\vu^T g_{\rm c}(\mL) - \vu^T g_{\rm c}({\overline \mL})_{:,j} \vone_n^T) = (\ve^{(i)})^T,
\end{equation*}
so $i$-th column of $g_{\rm c}(\mL)$ will turn to
\begin{equation*}
    g_{\rm c}(\mL)_{:,i} \mapsto
    g_{\rm c}(\mL)_{:,i} + ((\mA_{\overline \mL})_{:,j} - g_{\rm c}({\overline \mL})_{:,j})
    = g_{\rm c}(\mL)_{:,i} + ((\mA_\mL)_{:,i} - g_{\rm c}(\mL)_{:,i})
    = (\mA_\mL)_{:,i},
\end{equation*}
which is the desired output.
In conclusion, this layer maps the column $g_{\rm c}(\overline \mL)_{:,j}$ to $(\mA_{\overline \mL})_{:,j}$, without affecting any other columns.

As seen above, we need one layer per each unique value of $\vu^T g_{\rm c}(\mL)_{:,j}$ for each $\mL \in \wt \sG_{\delta}$. Note that there are $O(n(1/\delta)^{dn} / n!)$ such numbers, so we can use $O(n(1/\delta)^{dn} / n!)$ layers to finish our construction.


%%%%%%%%%%%%%%%%%%%%%%%%%%%%% 
%%%%%%%%%%%%%%%%%%%%%%%%%%%%% 

\section{Proof of Theorem~\ref{thm:anyseq2seq}}\label{sec:proof_anyseq2seq}
Proof of Theorem~\ref{thm:anyseq2seq} can be done in a similar way as Theorem~\ref{thm:univapprox}. As in the proof of Theorem~\ref{thm:univapprox}, there are three parts: Lemma~\ref{lemma:part1}, Proposition~\ref{prop:part2}, and Lemma~\ref{lemma:part3}. The statement and proof of Lemmas~\ref{lemma:part1} and \ref{lemma:part3} can be done in almost the same way, this time without permutation equivariance.

For the proof of the second part, which corresponds to Proposition~\ref{prop:part2}, we construct the network in a similar way.
Recall that we can assume without loss of generality that $\mX \in [0,1]^{d \times n}$. Choose
\begin{equation*}
    \mE = 
    \begin{bmatrix}
    0 & 1 & 2 & \cdots & n-1\\
    0 & 1 & 2 & \cdots & n-1\\
    \vdots & \vdots & \vdots & & \vdots\\
    0 & 1 & 2 & \cdots & n-1
    \end{bmatrix}.
\end{equation*}
Then, the first column of $\mX+\mE$ is in $[0,1]^d$, second is in $[1,2]^d$, and so on; this means that for all rows, the coordinates are monotonically increasing.
So we can use the same technique as the proof of Proposition~\ref{prop:part2} to divide the input values into cubes, quantize them to $\mL$, apply contextual mapping, and then value mapping. We describe each step in the following.

%%%%%%%%%%%%%%%%%%%%%%%%%%%%% 

\subsection{Quantization by feed-forward layers}
In a similar way as Lemma~\ref{lemma:quantize}, the goal of this step is to quantize the input in $[0,1]^d \times [1,2]^d \times \dots \times [n-1,n]^d$ to its discrete version:
\begin{equation*}
[0:\delta:1-\delta]^d \times [1:\delta:2-\delta]^d \times \dots \times [n-1:\delta:n-\delta]^d.
\end{equation*}

This can be done by $dn/\delta$ feed-forward layers.
We add $dn/\delta$ layers of the following form, for $k = 0, \delta, \dots, n-\delta$ and $i = 1, \dots, d$:
\begin{equation*}
    \mZ \mapsto \mZ + \ve^{(i)} \phi( (\ve^{(i)})^T \mZ - k\delta \vone_n^T),
    ~~
    \phi(t) = 
    \begin{cases}
    0 & t < 0 \text{ or } t \geq \delta\\
    -t & 0 \leq t < \delta.
    \end{cases}
\end{equation*}
After $dn/\delta$ layers, any input entry of $\mX + \mE$ in $[k\delta, k\delta+\delta)$ is quantized to $k\delta$.

%%%%%%%%%%%%%%%%%%%%%%%%%%%%% 

\subsection{Contextual mapping by attention layers}
By Step~1, we quantized any input $\mX + \mE$ to its quantized version. We call this quantized version $\mL$:
\begin{equation*}
\mL \in [0:\delta:1-\delta]^d \times [1:\delta:2-\delta]^d \times \dots \times [n-1:\delta:n-\delta]^d.
\end{equation*}
As done in Lemma~\ref{lemma:contextmap}, we define $\vu \defeq (1, \delta^{-1}, \dots, \delta^{-d+1})$ and $l_j \defeq \vu^T \mL_{:,j}$, for all $j \in [n]$.
Note that, because $\mL_{:,j} \in [j-1:\delta:j-\delta]^d$, we have
\begin{equation*}
    (j-1)(1+\delta^{-1}+\dots+\delta^{-d+1}) \leq l_j \leq (j-1)(1+\delta^{-1}+\dots+\delta^{-d+1}) + \delta^{-d+1} - \delta,
\end{equation*}
and $l_1 < l_2 < \dots < l_n$. Notice that this corresponds to the Category~1 in the proof of Lemma~\ref{lemma:contextmap}.

For simplicity of notation, let $s_j = (j-1) \sum_{k=0}^{d-1} \delta^{-k}$.
We stack $n(1/\delta)^d$ attention layers, with attention parts $\delta^{-d} \Psi(\cdot;l-\delta/2,l+\delta/2)$ for each $l \in \bigcup_{j=1}^n [s_j :\delta: s_j + \delta^{-d+1} - \delta]$, in increasing order of $l$.

These $n(1/\delta)^d$ attention layers perform selective shift operations on $l_j$'s, in increasing order of $j$. As seen in Appendix~\ref{sec:proof-category1}, shift operations result in $\wt l_1 < \wt l_2 <\dots<\wt l_n$. Also, the map from $\mL$ to $\wt l_n$ is one-to-one, which can be shown in the same way as Appendix~\ref{sec:proof-sublemma:1}. Since the range of $l_j$'s are a bit different, we have a different upper bound on $\wt l_n$:
\begin{align*}
    \wt l_n
    &\defeq l_n + \sum_{k=1}^{n-1} \delta^{-kd}(l_{n-k}-l_{n-k+1}) + \delta^{-nd} (l_n-l_1)\\
    &\leq l_n + \delta^{-d} (l_1 - l_n) + \delta^{-nd}(l_n - l_1)
    \leq s_n + \delta^{-d+1} - \delta + (\delta^{-nd} - \delta^{-d}) (s_n + \delta^{-d+1} - \delta)\\
    &=(\delta^{-nd} - \delta^{-d} + 1) \left ( (n-1) \frac{\delta^{-d}-1}{\delta^{-1}-1} + \delta^{-d+1} - \delta  \right )\\
    &\leq (\delta^{-nd} - \delta^{-d} + 1) (\delta^{-d} - 1) (n-1+\delta)
    < n\delta^{-(n+1)d}.
\end{align*}

Finally, we add an extra single-head attention layer with attention part $n\delta^{-(n+1)d-1} \psi(\cdot;0)$. We define the output of this layer as $g_{\rm c}(\mL)$. In a similar way as Appendix~\ref{sec:proof-category1}, this layer shifts all the layers by $n \delta^{-(n+1)d-1} \wt l_n$, thus making the intervals corresponding to different values of $\wt l_n$ disjoint from each other. This ensures that different contexts $\mL$ are mapped to distinct numbers in $\vu^T g_{\rm c}(\mL)$, thus implementing a contextual mapping.

%%%%%%%%%%%%%%%%%%%%%%%%%%%%% 

\subsection{Function value mapping by feed-forward layers}
Now, it is left to map $g_{\rm c}(\mL)$ to the desired output.
As seen in the last step, each different context $\mL$ maps to $n$ unique numbers $\vu^T g_{\rm c}(\mL)$, which are at least $\delta$ apart from each other. The value mapping step can be done in a similar way as Lemma~\ref{lemma:memorize}. 
The construction now requires $O(n(1/\delta)^{dn})$ layers because there is no permutation equivariance.


